\documentclass[reqno, 12pt]{amsart}

\usepackage{amssymb}
\usepackage{amsthm}
\usepackage{amsmath}
\usepackage{amsxtra}
\usepackage{mathrsfs}
\usepackage{comment}
\usepackage{graphicx}
\usepackage{placeins}
\usepackage{multicol}
\usepackage{bbm}
\usepackage{stmaryrd}
%\usepackage{enumerate}
\usepackage[inline, shortlabels]{enumitem}
\usepackage{mathdots}
\usepackage{colonequals}
\usepackage{hyperref}


\usepackage{calc}
\usepackage{tikz}
\usetikzlibrary{arrows,calc,automata,shadows,backgrounds,positioning,intersections,fadings,decorations.pathreplacing,shapes,matrix}
\usepackage{tikz-cd}

\usepackage{fullpage}

\newtheorem{thm}{Theorem}[section]
\newtheorem{lem}[thm]{Lemma}
\newtheorem{cor}[thm]{Corollary}
\newtheorem{conj}[thm]{Conjecture}
\newtheorem{prop}[thm]{Proposition}
\theoremstyle{definition}  
\newtheorem{defn} [thm] {Definition} 
\newtheorem{claim}[thm]{Claim}
\newtheorem{example} [thm] {Example}
\newtheorem{rem} [thm] {Remark}

\newenvironment{problab}[1]
{\noindent\textbf{Problem #1}.}
{\vskip 6pt}

\newenvironment{exolab}[1]
{\noindent\textbf{Exercise #1}.}
{\vskip 6pt}

\theoremstyle{remark}
\newtheorem*{soln}{Solution}

\newcommand{\C}{\mathbb C}
\renewcommand{\H}{\mathbb H}
\newcommand{\D}{\mathbb D}
\newcommand{\F}{\mathbb F}
\newcommand{\Q}{\mathbb Q}
\newcommand{\R}{\mathbb R}
\newcommand{\Z}{\mathbb Z}
\newcommand{\T}{\mathbb T}
\renewcommand{\P}{\mathbb P}
\renewcommand{\O}{\mathcal O}
\newcommand{\m}{\mathfrak m}
\newcommand{\A}{\mathbb A}
\renewcommand{\SS}{\mathbb S}

\newcommand{\B}{\mathcal B}
\newcommand{\E}{\mathcal E}

\newcommand{\SSpan}{\text{span}}

\DeclareMathOperator{\lcm}{lcm}
\DeclareMathOperator{\img}{img}
\DeclareMathOperator{\Ann}{Ann}
\DeclareMathOperator{\Tor}{Tor}
\DeclareMathOperator{\tr}{tr}
\DeclareMathOperator{\Hom}{Hom}
\DeclareMathOperator{\End}{End}
\DeclareMathOperator{\Aut}{Aut}
\DeclareMathOperator{\Gal}{Gal}
\DeclareMathOperator{\sgn}{sgn}
\DeclareMathOperator{\ord}{ord}
\DeclareMathOperator{\ad}{ad}
\DeclareMathOperator{\coker}{coker}
\DeclareMathOperator{\rank}{rank}

\DeclareMathOperator{\id}{id}

\usepackage{mathpazo}

\everymath{\displaystyle}

\tikzset{commutative diagrams/.cd, arrow style = tikz, diagrams = {>=latex}}

\newenvironment{amatrix}[1]{%
  \left(\begin{array}{@{}*{#1}{c}|c@{}}
}{%
  \end{array}\right)
}

\usepackage{abstract}
\renewcommand{\abstractname}{}    % clear the title
\renewcommand{\absnamepos}{empty} % originally center
%\setlength{\abstitleskip}

\begin{document}
%\renewcommand{\abstractname}{\vspace{-\baselineskip}}

\title{18.700 Problem Set 1}
%\author{Évariste Galois}
\dedicatory{Due Thursday, September 12 at 11:59 pm on \href{https://canvas.mit.edu/courses/27315}{Canvas}}
%\thanks{}
%\contrib{}
%\address{}

\maketitle

\begin{abstract}
  \noindent Collaborated with:\\
  Sources used: 
\end{abstract}
\vskip 0.5cm

\begin{problab}{1} (3 points)
  Do the planes $x_1 + 2 x_2 + x_3 = 4$, $x_2 - x_3 = 1$ and $x_1 + 3 x_2 = 0$ have at least one common point of intersection?
\end{problab}

\begin{problab}{2} (3 points)
  Give an example of a $3 \times 3$ matrix $A$ with real entries whose reduced row echelon form is
  \begin{equation*}
    \begin{pmatrix}
      1 & 0 & 1/2\\
      0 & 1 & 1/3\\
      0 & 0 & 0
    \end{pmatrix} \, ,
  \end{equation*}
and such that every entry of $A$ is a nonzero integer.
\end{problab}

\begin{problab}{3} (5 points)
  Fix $a,b,c,d \in \F$ and consider the matrix
  \begin{equation*}
    A \colonequals
    \begin{pmatrix}
      a & b & 1 & 0\\
      c & d & 0 & 1
    \end{pmatrix} \, .
  \end{equation*}
  Assuming $a \neq 0$ and $c \neq 0$, compute the reduced row echelon form of $A$. (\textit{Hint}: You will have to deal with two cases, depending on whether some quantity is zero or not.)
\end{problab}

\begin{problab}{4} (4 points)
  For each of the following augmented matrices, determine the value(s) of $h$ such that corresponding linear system is consistent.
  \begin{enumerate}[(a)]
    \item 
      \begin{equation*}
        \begin{amatrix}{2}
          1 & h & -3\\
          -2 & 4 & 6
        \end{amatrix}
      \end{equation*}
     \item 
      \begin{equation*}
        \begin{amatrix}{2}
          1 & 3 & -2\\
          -4 & h & 8
        \end{amatrix}
     \end{equation*}
  \end{enumerate}
\end{problab}

\begin{problab}{5} (3 points)
  Let
  \begin{equation*}
    S \colonequals
    \left\{
      \begin{pmatrix} 1\\ -2\\ 1\end{pmatrix} t +
      \begin{pmatrix} 2\\ 3\\ 0\end{pmatrix} : t \in \F
    \right\} \, .
  \end{equation*}
  Give a linear system whose solution set is $S$.
\end{problab}

\begin{problab}{6} (6 points)
  \begin{enumerate}[(a)]
    \item
      Suppose that $p(t) = a_0 + a_1 t + a_2 t^2$ is a quadratic polynomial with $a_0, a_1, a_2 \in \F$, whose graph passes through the points $(1,12), (2,15)$, and $(3,16)$. Find the coefficients $a_0, a_1, a_2$ by solving the following linear system.
      \begin{equation*}
        \begin{array}{cccc}
          a_0 & + a_1(1) & + a_2(1)^2 &= 12\\
          a_0 & + a_1(2) & + a_2(2)^2 &= 15\\
          a_0 & + a_1(3) & + a_2(3)^2 &= 16\\
        \end{array}
      \end{equation*}
    \item
      Suppose that $p(t) = a_0 + a_1 t + \cdots + a_n t^n$ is a polynomial of degree $n$ with $a_0, a_1, \ldots, a_n \in \F$, whose graph passes through the points $(u_1, v_1), (u_2, v_2), \ldots, (u_{n+1}, v_{n+1})$. Determine a linear system that the coefficients $a_0, a_1, \ldots, a_n$ must satisfy, and write down its corresponding augmented matrix.
  \end{enumerate}
\end{problab}

\end{document} 	
